\documentclass[10pt,a4paper]{amsart}
\usepackage[margin=19mm]{geometry}
\usepackage{amsmath,amssymb,mathtools}
\usepackage[scr=rsfs,cal=euler]{mathalfa}
\usepackage{xfrac}
\usepackage[unicode,hidelinks,bookmarks=false]{hyperref}
\newcommand{\ie}{i.e.\ }
\newcommand{\cf}{cf.\ }
\newcommand{\resp}{respectively\ }
\newcommand{\Subsection}{\subsection}
\providecommand{\nobreakdash}{\nobreak\hbox{-}}
\setlength{\emergencystretch}{2em}
\multlinegap0pt
\usepackage{amssymb}
\usepackage{algorithm}\usepackage{algpseudocode}
\newtheorem{prop}{Proposition}
\usepackage{cleveref}
\crefname{prop}{Proposition}{Propositions}
\title{On a Keller-Segel type equation to model Brain Microvascular Endothelial Cells growth's patterns}
\author{B. Ambrosio, A.F. Garroudji, S. Fitzsimons, H. Zaag, F.M. Elahi, E. Stratigakou-Polychronaki}
\date{Source: arXiv:2604.25180v2 (CC BY 4.0)}
\begin{document}
\maketitle

\noindent\emph{Source and licence.} On a Keller-Segel type equation to model Brain Microvascular Endothelial Cells growth's patterns. Version arXiv:2604.25180v2 (CC BY 4.0), distributed by arXiv under the Creative Commons Attribution 4.0 International licence, http://creativecommons.org/licenses/by/4.0/, as recorded in the arXiv licence metadata for this exact version and shown on the abstract page https://arxiv.org/abs/2604.25180v2. Full licence notice: \texttt{licenses/arxiv-2604.25180-CC-BY-4.0.txt}.\par\smallskip

\noindent\textit{Editorial scope.} Counted unit: the article's prop prop (source0.tex lines 544--546). The article's complete original proof of it is delivered with it (source0.tex lines 547--549, one \texttt{proof} environment of 3 lines). No mathematics is written, repaired or re-derived: the statement and the proof are the article's own lines, in the article's own notation, and no retained line is substituted or re-flowed.  Retained uncounted as the source-local context the proof needs: lines 526--536.  Nothing was omitted from any retained span, so the package carries no omission record.  No retained line contains a citation, so the wrapper carries no bibliography and no bibliography entry.  No retained line contains a figure environment, an image-inclusion command or drawing code, so no image is delivered and the package declares no asset dependency.  Editorial formatting only, in addition: an AMS \texttt{amsart} preamble that restates the article's own declarations and macros used by the retained lines, namely the environments \texttt{prop} on separate counters, together with the standard packages those lines need.  The retained environments print in this document as Proposition 1 in order of appearance, whereas the article prints the same items with the numbering of its own full document.  The complete native source of the article as distributed in the arXiv e-print for this exact version is bundled unchanged as \texttt{sources/199279/source0.tex}.
\begin{equation}
\label{eq:2times2}
\left \{
\begin{array}{rl}
      u_{1t}=&  f(u_1)-bu_1(v_2-v_1)+d_u(u_2-u_1) \\
      v_{1t}=& cu_1-ev_1+d_v(v_2-v_1)\\
         u_{2t}=&  f(u_2)-bu_2(v_1-v_2)+d_u(u_1-u_2) \\
      v_{2t}=& cu_2-ev_2+d_v(v_1-v_2)
 \end{array}
 \right.
\end{equation}

\begin{prop}
 The set $G=\{u_1>0,v_1>0,u_2>0,v_2>0\}$ is positively invariant for \Cref{eq:2times2}. 
\end{prop}

\begin{proof}
    We note that the set $u_1=u_2=0$ is an invariant manifold. Next, we observe that $u_1'>0$ at $u_1=0$ and $u_2>0$. The same argument holds for the other variables.   
\end{proof}

\end{document}
